\documentclass[11pt]{article}
\usepackage[T1]{fontenc}
\usepackage{lmodern,microtype,amsmath,amssymb,amsthm,mathtools}
\usepackage[margin=1in]{geometry}
\usepackage{tikz-cd,enumitem,booktabs,xurl}
\usepackage[hidelinks]{hyperref}
\hypersetup{pdftitle={Coalgebraic higher types and determinacy}}
\setlength{\emergencystretch}{2em}
\newtheorem{theorem}{Theorem}[section]
\newtheorem{proposition}[theorem]{Proposition}
\newtheorem{corollary}[theorem]{Corollary}
\theoremstyle{definition}
\newtheorem{definition}[theorem]{Definition}
\newcommand{\N}{\omega}
\newcommand{\B}{\omega^\omega}
\newcommand{\Pcat}{\mathcal P}
\newcommand{\Vcat}{\mathcal V}
\newcommand{\Hcat}{\mathcal H}
\newcommand{\Rcat}{\mathcal R}
\newcommand{\pArr}{\mathbin{\Rightarrow_{\!p}}}
\newcommand{\Coalg}{\operatorname{Coalg}}
\newcommand{\Par}{\operatorname{Par}}
\newcommand{\Tot}{\operatorname{Tot}}
\newcommand{\supp}{\operatorname{supp}}
\newcommand{\dom}{\operatorname{dom}}
\newcommand{\id}{\operatorname{id}}
\newcommand{\ev}{\operatorname{ev}}
\newcommand{\ran}{\operatorname{ran}}
\newcommand{\den}[1]{\mathopen{[\mkern-3mu[}#1\mathclose{]\mkern-3mu]}}
\title{Coalgebraic higher types and determinacy}
\author{}
\date{Version 1 --- 21 October 2026}
\begin{document}
\maketitle
\begin{abstract}
Partial function types on numerical partial equivalence relations admit
faithful coalgebraic realizations with countable carriers and binary
observations. We compare this realization with restriction splitting and
regular exact completion, and give an arithmetic-oracle bound for its
finite-type fragment. A range obstruction occurs already at the type of
partial numerical functions. These constructions and the obstruction
hold in ZF. Determinacy supplies a different extension: strategies give
Mealy or Moore reductions of stream predicates, and real determinacy
gives continuous representative lifts on invariant comeager domains for
quotients by countable group actions. These lifts have exact-domain
asynchronous coalgebraic realizations. We identify the composition
condition on their domains and a determinacy obstruction to complete
real-valued invariants of numerical oracle presentations.
\end{abstract}

\section{Computational types and the ambient set theory}

A recursion category can occur as a category of computations inside
a larger category of types. Restriction splitting supplies domains;
exact completion supplies certain quotients. Function types require
further quotients of programs. The question here is whether the
resulting operations admit simultaneous coalgebraic realizations and
which parts of that construction depend on choice.

We use ZF as the base theory. Additional axioms are stated in the
results that use them. Write $\N=\{0,1,\ldots\}$ and
$2=\{0,1\}$. Fix the Cantor pairing
$\langle m,n\rangle=(m+n)(m+n+1)/2+n$ and its computable
projections. Fix an acceptable numbering $(\varphi_e)_{e\in\N}$
of the partial recursive numerical functions. Composition $gf$ means
first $f$, then $g$. Equality of partial maps includes equality of
domains. The restriction $\bar f$ is the partial identity on $\dom f$.

A partial equivalence relation (PER) is a symmetric transitive relation
$R\subseteq\N^2$. Put $\supp R=\{a:aRa\}$ and
$|R|=\supp R/R$. Let $\Pcat$ have PERs as objects and partial
functions $f:|R|\rightharpoonup|S|$ tracked by numerical programs
as arrows. A tracker $e$ must satisfy, for every $a\in\supp R$,
\[
 f([a]_R)\downarrow\ \Longleftrightarrow\ \varphi_e(a)\downarrow,
 \qquad
 f([a]_R)=[\varphi_e(a)]_S\quad\hbox{when defined}.
\]
In particular, convergence is invariant under $R$. Products use
numerical pairing, and a singleton PER gives a restriction terminal
object $\mathbf1$. Program composition and strict pairing make
$\Pcat$ a cartesian restriction category. Put
$\Vcat=\Tot(\Pcat)$, and let $\mathbf N$ be equality on $\N$.
The full subcategory on positive finite powers of $\mathbf N$ is the
numerical recursion category $\Rcat$.

Define $R\pArr S$ by
\[
 e\,(R\pArr S)\,e'
 \quad\Longleftrightarrow\quad
 \forall aRa'\;
 \begin{cases}
 \varphi_e(a),\varphi_{e'}(a')\text{ both diverge},\quad\text{or}\\
 \varphi_e(a),\varphi_{e'}(a')\text{ converge to $S$-related values}.
 \end{cases}
\]
This is a PER whose classes are precisely the arrows
$R\rightharpoonup S$. Universal evaluation descends to
$\ev:(R\pArr S)\times R\rightharpoonup S$. The parameter
theorem gives total currying and the bijection
\begin{equation}\label{eq:curry}
 \Pcat(Q\times R,S)\cong\Vcat(Q,R\pArr S).
\end{equation}
Indeed, if $e$ tracks $h$, a total compiler sends $q$ to an index
$s(e,q)$ satisfying
$\varphi_{s(e,q)}(a)\simeq\varphi_e(\langle q,a\rangle)$.
Input equivalence for $h$ proves extensionality of this compiler.
Evaluation proves the two inverse identities in \eqref{eq:curry}.
The bijection is natural under total changes of $Q$.

The computation PER $LS$ identifies indices whose computations on
input $0$ both diverge or return $S$-related values. Suspension and
forcing give $\Pcat(R,S)\cong\Vcat(R,LS)$. The unit returns an
immediate computation, and multiplication runs an outer computation
and then the computation it returns. These explicit compilers define
a monad on $\Vcat$, with Kleisli category $\Pcat$.

Let $\Hcat$ have finite type expressions generated from
$\mathbf1,\mathbf N$ using products, $L$, and $\pArr$ as objects,
with all $\Pcat$-arrows between their PER interpretations.
Retaining the expressions as object labels makes changes of numbering
unambiguous. This includes both strict function
types and lazy types. For the latter, set $V_\iota=\mathbf N$,
$C_\sigma=LV_\sigma$, and
$V_{\sigma\to\tau}=C_\sigma\pArr V_\tau$.
A lazy function receives a computation name. Abstraction returns a
function value; application forces the function computation and passes
the argument name to it. The substitution theorem for these types is
proved in \cite[Section 6]{Application}.

All these definitions are available in ZF. Every class of a numerical
PER has a least member, so $|R|$ is canonically bijective with
$\{a\in\supp R:\forall b<a\;\neg(bRa)\}$. This selection is
set-theoretic; the evaluation and currying algorithms above do not
require it to be computable.

\section{Comparison with the existing completions}\label{sec:completion}

Let $K_r$ denote splitting of restriction idempotents. Let
$\Rcat^1$ have the positive finite powers of $\N$ and a singleton
as objects, with all partial recursive maps between them. The
singleton is its restriction terminal object. Set
$\mathcal K=K_r(\Rcat^1)$,
$\mathcal D=\Tot(\mathcal K)$, and
$\mathcal E=\operatorname{Ex}_{\mathrm{reg}}(\mathcal D)$.
Here $\mathcal D$ is equivalent to the regular category of c.e.
numerical supports and partial recursive maps total on their supports.
Its regular exact completion has the presentation by c.e. PERs and
c.e. saturated total functional relations described in
\cite[section ``Completions of categories of partial operators'']{External}.
Write $\Pcat_{\mathrm{ce}}$ for the full subcategory of $\Pcat$
on c.e. PERs.

\begin{proposition}\label{prop:completion}
There are fully faithful functors
$\mathcal K\to\Pcat$ and $\mathcal E\to\Vcat$.
The first identifies split domains with equality PERs on c.e. supports.
The second identifies $\mathcal E$ with the full subcategory of
$\Vcat$ on c.e. PERs. Moreover,
\[
 \Par(\mathcal E,\operatorname{Mono}(\mathcal E))
       \simeq\Pcat_{\mathrm{ce}}.
\]
\end{proposition}
\begin{proof}
A restriction idempotent on a numerical object is a partial identity
on a c.e. set. Its splitting therefore has precisely the indicated
support and partial recursive arrows, proving the first assertion.

A c.e. total functional relation $F$ between PERs $R,S$ is executed
by enumerating $F(a,b)$ until an output $b$ appears. Totality ensures
termination on valid inputs, and functionality ensures that all
outputs belong to the same $S$-class. Conversely a total tracker $e$
defines the saturated c.e. relation
\[
 F_e(a,b)\quad\Longleftrightarrow\quad
 aRa\ \land\ \exists c\,[\varphi_e(a)\downarrow=c\ \land\ cSb].
\]
Extensionality supplies saturation in the first coordinate. These
constructions identify the homsets, and represented functions determine
identities and composition.

For a partial tracker, the name-domain
$D_e=\{a:aRa\land\varphi_e(a)\downarrow\}$ is c.e. and
$R$-saturated. Restricting $R$ to $D_e$ gives a domain span in
$\mathcal E$. Conversely, consider a span
$R\xleftarrow{m}T\xrightarrow{g}S$ with $m$ monic. Its underlying
map on classes is injective: distinct classes are separated by maps
from $\mathbf1$. On input $a$, search for a valid $T$-name whose
$m$-image is $R$-related to $a$, and apply a tracker of $g$ to it.
The search is effective because the presentation is c.e. Monicity
makes the output class independent of the chosen preimage. Its
convergence domain is exactly the domain of the span. This proves
realizability of every span. Applied to the mono $m$, the same
search gives its inverse onto its c.e. saturated image in $R$.
Thus the span is isomorphic to the reconstructed domain span,
which also proves faithfulness of the partial-map comparison.
\end{proof}

This comparison does not identify all PERs with exact-completion
objects. The support of $\mathbf N\pArr\mathbf N$ is $\N$,
but its relation is extensional equality of partial recursive
functions and is not c.e. Its total-function subspace has the
additional non-c.e. support
$\mathrm{TOT}=\{e:\varphi_e\text{ is total}\}$.
The obstruction in \cite[section ``Total function spaces and finite
observations'']{External} shows that $\mathcal E$ has no exponential
of the numerical object by itself. The larger category supplies
operations unavailable in that completion.

\section{A simultaneous coalgebraic realization}\label{sec:realization}

Let $M(X)=2\times X^{\{+,-\}}$. An $M$-coalgebra is a set with
a binary output and two transitions. Partial coalgebra morphisms have
transition-closed subcoalgebra domains. Let $Q$ have carrier
$\mathbb Z$, output $1$ precisely at $0,1$, and transitions adding
or subtracting $1$. A coalgebra endomorphism of $Q$ commutes with
translations, so it is $n\mapsto n+k$. Preservation of the two
marked positions forces $k=0$. A nonempty subcoalgebra of $Q$ is
all of $Q$.

For a PER $R$, define $J(R)$ to be the coproduct of one copy of
$Q$ for each element of $|R|$:
\[
 J(R)=|R|\times\mathbb Z,\qquad
 o([a],n)=\chi_{\{0,1\}}(n),\qquad
 \delta_\pm([a],n)=([a],n\pm1).
\]
For $f:R\rightharpoonup S$, define
$Jf([a],n)=(f([a]),n)$ on $\dom f\times\mathbb Z$.
This uses the rigid marker construction of \cite[Section 4]{Catalogue}
with a single label.

\begin{theorem}\label{thm:coalgebra}
In ZF, $J:\Pcat\to\Par(\Coalg(M))$ is a faithful functor
with countable carriers. It preserves restrictions, zeros, and binary
near products. Its restriction to $\Hcat$ realizes application,
total currying, the computation monad, and lazy substitution
simultaneously in the category of represented arrows.
\end{theorem}
\begin{proof}
Every subcoalgebra of $J(R)$ is a union of whole copies of $Q$.
A map between two such unions sends each copy to one target copy,
and rigidity of $Q$ forces it to preserve integer phase. Thus the
construction $X\mapsto X\times\mathbb Z$, for arbitrary sets $X$,
embeds $\mathbf{Pfn}$ fully faithfully into partial coalgebra maps.
The functor $J$ is the faithful underlying-set
functor of $\Pcat$ followed by this marker embedding. This proves
functoriality, faithfulness, and preservation of restrictions and zeros.
Least representatives code each $|R|\times\mathbb Z$ into $\N$.

The coalgebra product of $J(R)$ and $J(S)$ is
$|R|\times|S|\times\mathbb Z$. Indeed, two maps from any
coalgebra into the two factors assign the same phase to each source
state: their binary outputs after all words in $+,-$ coincide, and
these outputs determine the phase uniquely. The two maps therefore
pair into the displayed coalgebra. The bijection
$|R\times S|\cong|R|\times|S|$ gives the canonical product
comparison, including the intersection of partial domains.

Every operation is an arrow of $\Pcat$, so $J$ sends its typed
equations to equalities of partial coalgebra morphisms. Total arrows
remain total. The represented category, with source objects retained
as tags, is isomorphic to $\Pcat$; hence \eqref{eq:curry} remains
its representing bijection. The monad on its total arrows is
$JR\mapsto J(LR)$, with unit and multiplication obtained by
applying $J$. These statements concern represented arrows. The
terminal marker $Q=J\mathbf1$ is terminal on total maps in the
marker subcategory, rather than the final coalgebra of all of
$\Coalg(M)$.
\end{proof}

For example, after the product comparison, application is the concrete
partial coalgebra morphism
\[
 (([e],[a]),n)\longmapsto([\varphi_e(a)],n).
\]
Its domain contains the entire integer orbit whenever the computation
converges. Currying sends $([q],n)$ to $([s(e,q)],n)$, without
executing the body. For lazy arguments, the $[a]$ coordinate is a
class in a computation PER; the function decides whether to force it.
Thus the construction retains both divergence and the multiplicity
of states having identical observations. The final Moore behavior of
$([a],n)$ depends only on $n$: the value $[a]$ is retained by its
orbit copy. Application is realized by a morphism between these
systems; the shifts themselves do not execute the application program.

\subsection{An effective bound and oracle relatives}

Write $0^{(\omega)}=\{\langle n,k\rangle:k\in0^{(n)}\}$
for the join of the finite Turing jumps. It has the degree of
arithmetical truth.

\begin{proposition}\label{prop:effective}
For finite type expressions defining $\Hcat$, the coalgebra carriers,
outputs, transitions, and represented arrow maps have presentations
uniformly computable relative to $0^{(\omega)}$. For the category
of partial $A$-recursive numerical maps, the same construction has
bound $A^{(\omega)}$.
\end{proposition}
\begin{proof}
Induction on a type expression produces an arithmetical formula for
its PER. Products use pairing and conjunction; $L$ and $\pArr$
add numerical quantifiers and halting predicates. Arithmetical truth
therefore decides support and equivalence uniformly across all finite
types. It also decides the canonical-carrier condition
$aRa\land\forall b<a\;\neg(bRa)$.

On a canonical input $a$, execute a supplied tracker $e$. If it
returns $b$, choose the least $c\leq b$ with $cSb$ using the
oracle. This computes the represented partial map with its exact
domain. Integer shifts and marker outputs require no further oracle.
All tests relativize to $A$. On the original numerical objects the
PERs are equality, so canonicalization is trivial and the arrow
action retains its original computability.
\end{proof}

Let $\varphi,\psi$ be two acceptable numberings, and let
$\Hcat_\varphi,\Hcat_\psi$ be their expression-labelled type
categories. Changing the numbering gives a concrete comparison of
their representations. The numerical translations are extended recursively
to products and computation types. Write $t_R,u_R$ for forward and
backward name translations at type $R$. At $R\pArr S$, compile
in the new numbering the program
$a\mapsto t_S(\varphi_e(u_R(a)))$. The parameter theorem produces
this code totally. Induction gives quotient bijections $b_R$
and preserves convergence. The equivalence
$F:\Hcat_\varphi\to\Hcat_\psi$ is the identity on type
expressions and sends $f:R\rightharpoonup S$ to
$b_Sfb_R^{-1}$. Its coalgebra comparison is
\[
 \kappa_R:J(R)\longrightarrow J'(FR),\qquad
 ([a]_R,n)\longmapsto([t_R(a)]_{FR},n)
\]
is a coalgebra isomorphism. For every partial $f:R\rightharpoonup S$,
the equality $J'(Ff)\kappa_R=\kappa_SJf$ includes equality of
domains. The same translations commute with evaluation, suspension,
and currying. Write $c^P_{R,S}:P(R\times S)\to PR\times PS$
for a functor's canonical product comparison. Then
\[
 c^{J'F}_{R,S}\kappa_{R\times S}
       =(\kappa_R\times\kappa_S)c^J_{R,S},
\]
since product names are translated componentwise. These are the
comparison equations of \cite[Section 3]{Catalogue}, applied to
typed cartesian restriction categories. Restriction to the numerical
recursion core gives a 1-cell in the representation 2-category
defined there.

\section{A range obstruction at the first function type}\label{sec:range}

A Heller range of $f:X\rightharpoonup Y$ is a restriction idempotent
$r$ with $rf=f$ such that $uf=vf$ implies $ur=vr$ for all
parallel $u,v:Y\rightharpoonup Z$. It is least among restriction
idempotents $d$ satisfying $df=f$: apply the condition to $d$ and
$\id_Y$. The following obstruction involves a generated type,
without adjoining an arbitrary additional PER.

\begin{theorem}\label{thm:no-range}
Let $S=\mathbf N\pArr\mathbf N$, and let
$z:\mathbf1\to S$ select the total-zero function. In $\Hcat$,
$z$ has no least supporting restriction idempotent and hence no
Heller range. Moreover, $z$ is not a partial isomorphism: there is
no $s:S\rightharpoonup\mathbf1$ with $sz=\id_{\mathbf1}$ and
$zs=\bar s$. These conclusions hold in ZF.
\end{theorem}
\begin{proof}
For each $n$, a program tracking $r_n:S\rightharpoonup S$
runs its input index $e$ on $n$, returning $e$ if the result is
$0$ and otherwise diverging. This is an extensional restriction
idempotent, and $r_nz=z$. If a least supporting restriction $r$
existed, then $r\leq r_n$ for every $n$. Together with $rz=z$,
this makes the numerical domain of every tracker of $r$ exactly
\[
 Z=\{e:\forall n\;\varphi_e(n)\downarrow=0\}.
\]
This set is not c.e. Let $K=\{k:\varphi_k(k)\downarrow\}$.
Uniformly in $k$, compile a program which, on input $n$, returns
$0$ if $\varphi_k(k)$ has not halted within $n$ steps and diverges
otherwise. Its index belongs to $Z$ exactly when $k\notin K$.
Thus a tracker of $r$ would enumerate the non-c.e. complement of $K$.

A partial isomorphism inverse $s:S\rightharpoonup\mathbf1$ would satisfy
$sz=\id_{\mathbf1}$ and $zs=\bar s$. Its domain would be precisely
the zero-function class, giving the same contradiction.
\end{proof}

The theorem also rules out an exact realization of the singleton
image by a tracked partial identity. It persists after adding DC,
AD, or real determinacy: its proof uses only numerical computation
and ZF.

\begin{proposition}\label{prop:existing-ranges}
The functor $J$ preserves every Heller range that exists in the full
PER category $\Pcat$.
\end{proposition}
\begin{proof}
Let $r$ be a Heller range of $f:R\rightharpoonup S$ and put
$I=\operatorname{im}f\subseteq|S|$. If $I$ is empty, applying
the range condition to $\id_S$ and $0$ gives $r=0$.
Otherwise choose a name $b$ with $[b]_S\in I$, set
$B=\{a\in\supp S:[a]_S\in I\}$, and define the PER
$T=S\cup(B\times B)$. The identity-name map $u:S\to T$
and the constant-$b$ map $v:S\to T$ are total tracked maps.
They satisfy $uf=vf$, so $ur=vr$. Since $r$ is a partial
identity, the latter equation forces $\dom r\subseteq I$;
$rf=f$ gives the reverse inclusion. Thus $r$ is the actual
image restriction, and $Jr$ selects exactly the image orbits of
$Jf$. That partial identity is its ambient Heller range.
\end{proof}

The ranges of the original recursion category are also
preserved by $J$. Dovetailing numerical computations until a preimage
is found gives, for each $f$, a partial section $s$ with
$fs=\ran(f)$. In any restriction
category, $r=fs$ and $rf=f$ imply
$uf=vf\Rightarrow ur=ufs=vfs=vr$. A restriction functor
preserves these equations. More generally, c.e. PER targets admit
the c.e. saturated name-image of each arrow from a c.e. PER, so the
comparison in Section~\ref{sec:completion} preserves their exact
image ranges. The obstruction concerns the additional function type.

\section{Dependent choice and countable coalgebras}

Dependent choice (DC) states that every serial relation on a nonempty
set admits a sequence following that relation from any specified
starting point. A relation is serial when every element has a successor.
Let $\mathcal P_+(X)$ denote the set of nonempty subsets of $X$.

\begin{proposition}
Over ZF, DC is equivalent to the assertion that every pointed
coalgebra $c:X\to\mathcal P_+(X)$ has an infinite run
$x_0,x_1,\ldots$ satisfying $x_{n+1}\in c(x_n)$.
\end{proposition}
\begin{proof}
Associate to $c$ the serial relation $xRy\iff y\in c(x)$.
Conversely, a serial relation defines this coalgebra by taking
successor sets. The two assertions are identical under this
correspondence.
\end{proof}

Deterministic transitions need no choice: their runs follow by
recursion. On a carrier with a specified injection into $\N$, even
nonempty successor sets have a least coded successor. DC becomes
relevant for unrestricted state spaces or successively chosen lifts.

Countable Boolean algebras also admit the ultrafilters needed in
countable germ representations without a choice axiom. For an
enumeration $b_0,b_1,\ldots$ of a Boolean algebra and a proper ideal
$I$, put $e_0=1$ and choose $e_{n+1}=e_n\wedge b_n$ if that
meet is outside $I$; otherwise put
$e_{n+1}=e_n\wedge\neg b_n$. The second meet is then outside
$I$, since the join of the two meets is $e_n$. The set
$U=\{b:\exists n\;(e_n\leq b)\}$ is an ultrafilter disjoint
from $I$. Each decision is specified by a formula, so this is a
ZF recursion. Applied to the proper ideals of agreement domains
for unequal arrows, it gives the separating labels in the countable
germ construction of \cite[Section 4]{Catalogue}. Coding a germ by
a label index and an arrow index avoids a countable choice of fibre
enumerations. This argument is restricted to explicitly countable
presentations.

The Boolean Prime Ideal Theorem (BPI), in its ultrafilter formulation,
asserts that every nontrivial Boolean algebra has an ultrafilter.
It extends this separation argument to arbitrary Boolean algebras:
apply it to the quotient by a proper ideal and pull back the
ultrafilter. Applied to $\mathcal P(\N)/\mathrm{Fin}$, it yields
a nonprincipal ultrafilter on $\N$. Under AD every ultrafilter is
countably complete \cite[Theorem 2.8]{KoellnerWoodin}; such an
ultrafilter on $\N$ would contain every cofinite tail and therefore
their empty intersection. Thus BPI is incompatible with AD, while
the countable separation argument above remains available in ZF.

\section{Determinacy as a coalgebraic comparison theorem}\label{sec:games}

Baire space $\B$ has the product topology of discrete $\N$.
For a finite string $s$, let $[s]$ be its basic cylinder. A causal
map $f:\B\to\B$ has $f(x)(n)$ determined by
$x\upharpoonright(n+1)$; a strictly causal map $g$ has
$g(y)(n)$ determined by $y\upharpoonright n$.
Causal maps have countable-state Mealy presentations, with output
and transition map $X\times\N\to\N\times X$. Strictly causal
maps have Moore presentations $X\to\N\times X^\N$.
Each presentation includes a specified initial state. Finite input
histories give the states in both constructions, with the empty
history as initial state.
Their tables can be coded as subsets of $\N$.

AD states that every length-$\omega$ perfect-information game with
integer moves is determined. A winning strategy is a function on
finite play histories, so it has a real code. PD restricts this
assertion to projective payoff sets; here projective means obtained
from Borel sets by finitely many projections and complements, with
real parameters allowed.

\begin{theorem}\label{thm:mealy-moore}
Assume ZF+AD. For $A,B\subseteq\B$, exactly one of the
following alternatives holds:
\begin{enumerate}[label=(\roman*),nosep]
\item a real-coded causal Mealy machine realizes $f$ with
      $A=f^{-1}(B)$;
\item a real-coded strictly causal Moore machine realizes $g$ with
      $\B\setminus B=g^{-1}(A)$.
\end{enumerate}
For projective $A,B$, PD suffices in place of AD.
\end{theorem}
\begin{proof}
Let I and II alternately produce $x(n)$ and $y(n)$, with II
winning exactly when $x\in A\iff y\in B$. A winning II
strategy produces $y(n)$ after seeing $x(n)$; its finite histories
give the Mealy machine in (i). A winning I strategy produces
$x(n)$ before seeing $y(n)$; its histories give the Moore machine
in (ii). Determinacy supplies one of these strategies. For
projective $A,B$, the payoff is projective.

Both alternatives cannot hold. Their composite $h=gf$ is strictly
causal and has a unique fixed stream. Construct its $n$th coordinate
from the preceding coordinates, filling later coordinates temporarily
with zero; strict causality makes the choice of that tail irrelevant.
Induction proves existence and uniqueness of the fixed stream $x$.
But then
$x\in A\iff f(x)\in B\iff gf(x)\notin A$, a contradiction.
\end{proof}

This is the strategy proof of Wadge's lemma, with the two timing
conventions made explicit \cite[Theorem 7D.3]{Moschovakis}.
The fixed stream is also obtained by feeding each Moore output back
as the next input. For a Moore presentation $(o,\delta)$, its
closed transition is $q\mapsto\delta_{o(q)}(q)$.
For Borel payoff sets the game is already determined by Borel
determinacy; AD extends the assertion to arbitrary stream predicates.
The familiar Wadge order and its well-foundedness under determinacy
are discussed in \cite{Carroy}.

The real code of a strategy is an oracle for its machine, not a
recursive program in general. A game requiring II to produce a fixed
nonrecursive stream has a winning strategy and no recursive winning
strategy. Thus determinacy gives additional real-coded comparisons
without changing which numerical programs halt.

\section{Representative lifts on invariant domains}\label{sec:lifts}

Real determinacy, written $\mathrm{AD}_{\mathbb R}$, permits a
whole real as each move. It implies AD. We use it with DC.
A subset of a Polish space has the Baire property if it differs
from an open set by a meager set; a set is meager if it is a
countable union of nowhere dense sets, and comeager if its complement
is meager. AD gives the Baire property to every set of reals
\cite[Theorem 2.7]{KoellnerWoodin}. A map into Baire space whose
basic-open preimages have this property is continuous on a comeager
subset \cite[p.~50]{Ikegami}; the proof below chooses that subset
to be invariant and $G_\delta$.

\begin{theorem}\label{thm:generic-lift}
Assume ZF+DC+$\mathrm{AD}_{\mathbb R}$. Let a specified countable
group $\Gamma$ act on $X=\B$ by homeomorphisms, and let $E$ be
its orbit equivalence relation. Let $S$ be an equivalence relation
on a nonempty $Y\subseteq\B$. Every function
$F:X/E\to Y/S$ has a representative lift
$f:D\to Y$ on an $E$-invariant comeager $G_\delta$ set
$D\subseteq X$, with $f$ continuous and
\begin{equation}\label{eq:quotient-lift}
 [f(x)]_S=F([x]_E)\qquad(x\in D).
\end{equation}
If instead ZF+DC+PD holds and the relation
$P(x,y)\iff y\in Y\land F([x]_E)=[y]_S$ is projective,
the same conclusion holds.
\end{theorem}
\begin{proof}
Every section of $P$ is nonempty. In the real-move game where I
first chooses $x$ and II then chooses $y$, let II win iff $P(x,y)$;
later moves are irrelevant. I has no winning strategy, since its
first move admits a suitable reply. A winning II strategy therefore
uniformizes $P$ by a total function $f_0:X\to Y$.

For each basic open set $U_n\subseteq\B$, choose an open
$V_n\subseteq X$ and a meager error set witnessing the Baire
property of $f_0^{-1}(U_n)$. DC permits these countably many
choices and their closed nowhere dense covers. Outside the union
of all errors, enlarged to a meager $F_\sigma$ set, the resulting
comeager $G_\delta$ set $G$ satisfies
$f_0^{-1}(U_n)\cap G=V_n\cap G$ for every $n$.
Hence $f_0|G$ is continuous.

Put $D=\bigcap_{\gamma\in\Gamma}\gamma^{-1}G$.
Homeomorphisms preserve comeagerness, so the Baire category theorem
makes $D$ a dense comeager $G_\delta$. The group law makes it
invariant. Restricting $f_0$ to $D$ proves \eqref{eq:quotient-lift}.

Under PD, projective uniformization supplies a projective graph for
$f_0$ \cite[Theorem 6C.5]{Moschovakis}. Its basic-open preimages
are projective and have the Baire property. The remaining argument
is unchanged.
\end{proof}

Writing $q_D:D\to D/E$ and $q_Y:Y\to Y/S$ for the quotient
maps, the equality just proved is the commutative square
\[
\begin{tikzcd}[column sep=large]
D\arrow[r,"f"]\arrow[d,"q_D"']&Y\arrow[d,"q_Y"]\\
D/E\arrow[r,"F|_{D/E}"']&Y/S.
\end{tikzcd}
\]
The lift preserves output classes; it need not choose the same
output name for equivalent inputs. Invariant uniformization requires
that stronger equality of names, a distinction studied in
\cite{KechrisWolman}.

\subsection{An exact-domain machine for the lift}

A pointed asynchronous transducer consists of a state set $Q$, an
initial state $q_0\in Q$, and a total transition and output rule
$Q\times\N\to(\N\sqcup\{\ast\})\times Q$, where $\ast$
means no output on that step. It defines a partial map on streams:
the domain consists of inputs producing infinitely many numerical
outputs, and the output stream omits the $\ast$ symbols.

\begin{proposition}\label{prop:transducer}
Every continuous $f:D\to\B$ with $D\subseteq\B$ a dense
$G_\delta$ has a real-coded countable-state asynchronous transducer
whose partial stream function is exactly $f$, including its domain.
\end{proposition}
\begin{proof}
Write $D=\bigcap_i O_i$ with each $O_i$ open. A state records
the input prefix and the number $n$ of outputs already emitted.
The initial state is the empty prefix with $n=0$.
After extending the input prefix to $s$, emit $m$ precisely when
\begin{equation}\label{eq:emission}
 [s]\subseteq\bigcap_{i\leq n}O_i
 \quad\hbox{and}\quad
 \forall x\in[s]\cap D\; f(x)(n)=m.
\end{equation}
Otherwise emit $\ast$. Density of $D$ ensures uniqueness of $m$.
Advance $n$ exactly when a numerical output is emitted.

For $x\in D$, continuity and openness imply that a sufficiently
long prefix satisfies \eqref{eq:emission} for each successive $n$.
Thus every output digit is eventually emitted and equals $f(x)(n)$.
If $x\notin D$, choose $j$ with $x\notin O_j$. No prefix
of $x$ lies inside $O_j$, so no output numbered $j$ or later
can be emitted. The productivity domain is exactly $D$.
The rule is a function on a countable set of states and inputs;
its numerical graph is a real code for the machine.
\end{proof}

Applied to Theorem~\ref{thm:generic-lift}, this gives a coalgebraic
realizer of the quotient function on the saturated domain $D/E$.
The transition table supplies an oracle for its execution. Because
its productivity domain is $G_\delta$, its name-level partial map
is an arrow of the continuous-operator recursion category
$\mathcal B_o$ of Montagna--Sorbi
\cite[Lemma 2.5, Definition 2.2, and Proposition 3.1]{MontagnaSorbi}.

\subsection{Composition of the obtained partial maps}

Fix Polish ambient spaces $X_0,X_1,X_2$. Regard the obtained maps
as partial continuous maps $f:D\subseteq X_0\to X_1$ and
$g:E\subseteq X_1\to X_2$, using the inclusion $Y\subseteq\B$
for a lift's codomain. Their composite has domain
$D\cap f^{-1}(E)$. It is again $G_\delta$ when $D,E$ are
$G_\delta$ in their respective ambient spaces. Quotient compatibility
and saturation also compose. Comeagerness does not: take $f$ constant
with value $y_0\in X_1=\B$ and
$E=\B\setminus\{y_0\}$. Both original domains can be
comeager while the composite has empty domain.

A sufficient condition is that $f$ preserve category in the
preimage sense: the inverse image of every meager subset of $X_1$
is meager in $D$. With $D$ comeager in $X_0$, this makes
$D\cap f^{-1}(E)$ comeager in $X_0$ whenever $E$ is comeager
in $X_1$. Such partial maps with comeager $G_\delta$ domains
are closed under composition. The lifting theorem does not impose
this condition, so its general conclusion belongs to the category
of partial maps with actual $G_\delta$ domains.

\section{Consequences for classification by integer sets}\label{sec:classification}

There is a distinction between giving each presentation a real code
and assigning a complete real-valued invariant to its equivalence
class. Determinacy imposes restrictions on the latter.

Under AD there is no transversal for eventual equality $E_0$ on
$2^\N$. To see this, suppose $T$ chose one point per class.
Its translates by finite coordinate flips are pairwise disjoint and
cover Cantor space. All have the Baire property. If $T$ is meager,
this is a countable meager cover. If $T$ is nonmeager, it is
comeager in some cylinder. A nontrivial flip beyond that cylinder's
specified coordinates makes a disjoint translate comeager in the
same cylinder, again a contradiction. Thus quotient sets of reals
behave differently from quotients of subsets of $\N$, whose least
representatives exist in ZF.

Let $\mathcal D_T$ be the Turing degrees of subsets of $\N$.
A cone is $\{d:d\geq_T b\}$ for some degree $b$. Under AD,
the sets of degrees containing a cone form a countably complete
ultrafilter \cite[Theorems 2.8--2.9]{KoellnerWoodin}.

\begin{proposition}\label{prop:cone-invariant}
Assume ZF+AD. Every Turing-degree-invariant function
$I:2^\N\to2^\N$ is constant on a Turing cone.
\end{proposition}
\begin{proof}
For each $n$, the cone ultrafilter chooses one of the two sets of
degrees on which $I(A)(n)$ is $0$ or $1$. Countable completeness
makes the intersection of these chosen sets contain a cone. All
coordinates of $I$ are fixed on that cone.
\end{proof}

Let $\Rcat_A$ have the positive finite powers of $\N$ as objects
and partial $A$-recursive functions as arrows. Then
$\Rcat_A=\Rcat_B$ exactly when $A\equiv_T B$:
simulation proves one direction, and the characteristic functions
of the oracles prove the other. Equality here is equality as concrete
subcategories of $\mathbf{Pfn}$. There is therefore no complete
real-valued invariant for equality of these presentations under AD.
Indeed, a cone contains both a degree and its strictly higher jump,
while Proposition~\ref{prop:cone-invariant} assigns them the same
invariant. For classification under abstract or structured equivalence,
the same obstruction applies if the corresponding categories do not
all become equivalent on a cone. Whether they do is a separate
categorical question.

\section{The resulting extension problem}

The numerical exact-completion comparison and the coalgebraic
realization are established in Sections~\ref{sec:completion} and
\ref{sec:realization}. They give higher-type operations and their
comparison maps without choice. Theorem~\ref{thm:no-range} prevents
this particular program-quotient extension from having ranges for
all its arrows. A completion retaining those ranges must change
the admitted arrows or the interpretation of function types.

Determinacy supplies real-coded reductions of behavioral predicates.
Real determinacy, or projective determinacy with projective data,
also supplies representative lifts realized by partial continuous
transducers. A higher-type construction based on those lifts must
specify composable domains and compatible choices of realizers.
Neither faithfulness of a coalgebra representation nor existence of
individual lifts alone asserts a universal property for that
construction.

\begin{thebibliography}{9}
\small
\bibitem{Application} \emph{Application and substitution for external
higher types}, version 1, 2026.
\bibitem{Catalogue} \emph{Recursion categories and their coalgebraic
representations}, version 3, 2026.
\url{https://github.com/flengyel/notes/tree/main/recursion_categories}.
\bibitem{External} F.~Lengyel, \emph{An External Finite-Type Completion
of Recursion Categories}, version 0.11, 2026.
\url{https://github.com/flengyel/External_Finite_Type_Completion}.
\bibitem{Moschovakis} Y.~N. Moschovakis, \emph{Descriptive Set Theory},
second edition, 2009, Theorems 6C.5 and 7D.3.
\url{https://www.math.ucla.edu/~ynm/lectures/dst2009/dst2009.pdf}.
\bibitem{Carroy} R.~Carroy, L.~Motto Ros, and S.~Scamperti,
\emph{A classification of the Wadge hierarchies on zero-dimensional
Polish spaces}, arXiv:2302.01319v2, 2023.
\url{https://arxiv.org/abs/2302.01319}.
\bibitem{Ikegami} D.~Ikegami, \emph{Real determinacy and real Blackwell
determinacy}, RIMS K\^oky\^uroku \textbf{1754} (2011), 32--74.
\url{https://www.kurims.kyoto-u.ac.jp/~kyodo/kokyuroku/contents/pdf/1754-05.pdf}.
\bibitem{KoellnerWoodin} P.~Koellner and W.~H. Woodin,
\emph{Large Cardinals from Determinacy}, author manuscript, Section 2.1.
\url{https://bpb-us-e1.wpmucdn.com/websites.harvard.edu/dist/f/94/files/2022/07/LCFD.pdf}.
\bibitem{KechrisWolman} A.~S. Kechris and M.~Wolman,
\emph{Invariant uniformization}, arXiv:2405.15111v3, 2026.
\url{https://arxiv.org/abs/2405.15111}.
\bibitem{MontagnaSorbi} F.~Montagna and A.~Sorbi,
\emph{Creativeness and completeness in recursion categories of partial
recursive operators},
Journal of Symbolic Logic \textbf{54} (1989), 1023--1041.
\href{https://doi.org/10.2307/2274763}{doi:10.2307/2274763}.
\end{thebibliography}
\end{document}
